\documentclass[12pt,a4paper,oneside]{book}

\usepackage[T1]{fontenc}
\usepackage{lmodern}
\usepackage[margin=2.5cm]{geometry}
\usepackage{amsmath}
\usepackage{amssymb}
\usepackage{graphicx}
\usepackage{fancyhdr}
\usepackage{hyperref}
\usepackage{setspace}
\usepackage{array}
\usepackage{booktabs}
\usepackage{multirow}
\usepackage{float}
\usepackage{caption}
\usepackage{subcaption}


% Python code highlighting
    language=Python,
    basicstyle=\ttfamily\small,
    breaklines=true,
    keywordstyle=\color{blue},
    commentstyle=\color{gray},
    stringstyle=\color{red},
    showstringspaces=false,
    frame=single,
    rulecolor=\color{black!30}
}

% Document metadata
\title{\Large \textbf{ECG Simulator}\\
        \large Comprehensive Technical Documentation\\
        \normalsize Phase 1/2 (Conduction Graph) and Phase 3 (FHN Cable)}
\author{Enrico Caruso\\Engineering \& Simulation}
\date{\today}

% Header and footer
\pagestyle{fancy}
\fancyhf{}
\fancyhead[L]{\small ECG Simulator Documentation}
\fancyhead[R]{\small Part 1 of 3}
\fancyfoot[C]{\thepage}

% Custom commands
\newcommand{\class}[1]{\texttt{\textbf{#1}}}
\newcommand{\func}[1]{\texttt{#1}()}
\newcommand{\param}[1]{\textit{#1}}

\onehalfspacing

\begin{document}

\frontmatter

\maketitle

\tableofcontents
\listoffigures
\listoftables

\chapter*{About This Documentation}

This document provides exhaustive technical documentation of the ECG Simulator project, designed for four distinct audiences simultaneously:

\begin{itemize}
    \item \textbf{Physicians and medical students}: ECG theory and clinical interpretation with minimal mathematical background required.
    \item \textbf{Bioscientists and physiologists}: Cardiac electrophysiology without requiring programming expertise.
    \item \textbf{Mathematicians and physicists}: Full mathematical rigor and physical modeling details.
    \item \textbf{Software engineers}: Complete implementation details, code structure, and dependencies.
\end{itemize}

\noindent All concepts are explained from first principles. Where appropriate, concepts are stratified so that readers can engage at their level of interest.

\mainmatter

% ================================================================
\chapter{Project Overview}
% ================================================================

\section{Purpose and Scope}

The ECG Simulator is a Python-based biomedical simulation framework that models the electrical activity of the human heart and generates realistic 12-lead electrocardiograms (ECGs) under normal and pathological conditions. The simulator serves as both a:

\begin{enumerate}
    \item \textbf{Educational tool}: Teaching ECG interpretation and understanding cardiac electrophysiology
    \item \textbf{Research platform}: Studying the relationship between conduction abnormalities and ECG morphology
    \item \textbf{Clinical decision support}: Generating reference ECGs for various pathological states
\end{enumerate}

\section{Key Capabilities}

The simulator can model the following clinical scenarios:

\begin{itemize}
    \item Normal sinus rhythm (NSR) at various heart rates
    \item Atrial pathologies: atrial fibrillation, atrial flutter, atrial premature contractions
    \item Atrioventricular (AV) conduction abnormalities: first-, second-, and third-degree AV blocks
    \item Bundle branch blocks: right bundle branch block (RBBB), left bundle branch block (LBBB)
    \item Ventricular arrhythmias: premature ventricular contractions (PVCs), ventricular tachycardia
    \item Accessory pathway conduction: Wolff-Parkinson-White (WPW) syndrome
    \item Sinus node dysfunction: sinus bradycardia, sinus tachycardia
\end{itemize}

\section{Development Phases}

The project is organized into three phases:

\subsection{Phase 1/2: Conduction Graph Engine}

The Phase 1/2 implementation uses a \textbf{discrete conduction graph} model where:

\begin{itemize}
    \item The heart is represented as a directed graph of \textbf{nodes} (anatomical regions: SA node, RA, LA, AV node, His bundle, ventricular segments)
    \item \textbf{Activation propagates} as breadth-first search (BFS) with physiological conduction velocities
    \item Each node has fixed \textbf{activation delay} from the previous node
    \item Refractory periods are \textbf{hardcoded} for each node
    \item \textbf{ECG dipole contributions} from each node are modeled as Gaussian functions in time
\end{itemize}

\noindent \textbf{Advantages}:
\begin{itemize}
    \item Extremely fast (microseconds per step)
    \item All pathologies implemented as parameter overrides
    \item Clinically calibrated timing and morphology
\end{itemize}

\noindent \textbf{Limitations}:
\begin{itemize}
    \item Refractory period is an input, not emergent from dynamics
    \item No continuous state evolution or tissue coupling
\end{itemize}

\subsection{Phase 3: FitzHugh-Nagumo Cable Engine}

The Phase 3 implementation uses a \textbf{1-D monodomain FitzHugh-Nagumo cable} model where:

\begin{itemize}
    \item The conduction system is represented as a \textbf{linear chain of cells} (18 total)
    \item Each cell solves the \textbf{FitzHugh-Nagumo ordinary differential equations} (ODEs)
    \item Action potential (AP) propagates via \textbf{diffusion coupling} between adjacent cells
    \item Refractory period \textbf{emerges} naturally from the slow recovery variable dynamics
    \item \textbf{Multi-beat rhythm} requires no explicit pacemaker programming
    \item ECG dipoles updated \textit{continuously} as the AP wavefront propagates
\end{itemize}

\noindent \textbf{Advantages}:
\begin{itemize}
    \item Emergent refractory period (no hardcoding)
    \item Continuous state evolution (more physiologically faithful)
    \item Can model effects of rate-dependent APD changes
\end{itemize}

\noindent \textbf{Limitations}:
\begin{itemize}
    \item Slower than Phase 1/2 (2.2$\times$ real-time on modern hardware)
    \item Pathologies less developed (manual parameter override)
    \item Timing requires careful calibration
\end{itemize}

\section{User-Selectable Engine}

The user can switch between \phaseI and \phaseIII engines via a toolbar combo box. The switch:
\begin{itemize}
    \item Is seamless (pause current engine, swap, resume)
    \item Reinitializes the new engine with current parameter values
    \item Does not modify parameter settings
\end{itemize}

% ================================================================
\chapter{Theoretical Foundations: Mathematical and Physical Models}
% ================================================================

\section{Cardiac Electrophysiology Primer}

\subsection{The Action Potential}

Every cardiac myocyte is electrically excitable. When stimulated, it generates an \textbf{action potential} (AP)---a stereotypical time-evolution of transmembrane voltage $V$ that drives heart contraction.

\subsubsection{Phases of the Cardiac Action Potential}

The AP has five distinct phases (Figure~\ref{fig:ap_phases}):

\begin{figure}[H]
    \centering
    
    \begin{center}
    \texttt{[Diagram - text description in document]}
    \end{center}
    
    \caption{Five phases of the cardiac action potential. Phase 0: rapid depolarization (inward $\text{Na}^+$ current). Phase 1: early repolarization (outward $\text{K}^+$ current). Phase 2: plateau (inward $\text{Ca}^{2+}$ current balances outward $\text{K}^+$). Phase 3: rapid repolarization. Phase 4: resting state.}
    \label{fig:ap_phases}
\end{figure}

\noindent \textbf{Phase 0 (Depolarization)}: The membrane potential $V$ rapidly depolarizes from resting $V_{\text{rest}} \approx -90$ mV to peak $\approx +30$ mV due to rapid inward $\text{Na}^+$ current. This phase lasts $\approx 1$ ms.

\noindent \textbf{Phase 1 (Early Repolarization)}: Brief outward $\text{K}^+$ current causes $V$ to decline slightly.

\noindent \textbf{Phase 2 (Plateau)}: Inward $\text{Ca}^{2+}$ current balances outward $\text{K}^+$ current, maintaining $V \approx 0$ mV. This phase lasts $100$--$300$ ms depending on cell type and conditions.

\noindent \textbf{Phase 3 (Late Repolarization)}: Dominant outward $\text{K}^+$ current repolarizes the cell back toward $V_{\text{rest}}$.

\noindent \textbf{Phase 4 (Resting/Diastole)}: The cell remains at rest potential, ready for the next stimulus.

\subsubsection{Refractoriness}

During and immediately after an AP, the cell is \textbf{refractory}---it cannot be re-excited by an incoming stimulus:

\begin{itemize}
    \item \textbf{Absolute refractory period}: During phases 0--2 and early phase 3, $\text{Na}^+$ channels are inactivated; no re-stimulation possible.
    \item \textbf{Relative refractory period}: Late phase 3, when $V$ is still below threshold. Re-stimulation is possible only with a \textbf{supernormal stimulus} (larger than normal).
\end{itemize}

The duration of refractoriness is roughly proportional to the AP duration (\textbf{APD}). Longer APD $\Rightarrow$ longer refractoriness.

\subsection{Conduction Velocity}

Electrical activation spreads from cell to cell via gap junctions (protein channels connecting adjacent myocytes). The speed of this propagation is the \textbf{conduction velocity} (CV):

$$\text{CV} \approx 0.5 \sqrt{\frac{D}{\epsilon}} \cdot \frac{\text{dx}}{\text{cell time constant}}$$

\noindent where:
\begin{itemize}
    \item $D$ is the diffusion coefficient (cytoplasmic resistivity)
    \item $\epsilon$ is the excitability (inversely related to refractory period)
    \item $\text{dx}$ is the cell size
\end{itemize}

Typical values in the heart:
\begin{itemize}
    \item Atrial myocardium: $0.5$--$1.0$ m/s
    \item His-Purkinje system: $2$--$4$ m/s (high CV)
    \item Ventricular myocardium: $0.3$--$0.5$ m/s
    \item AV node: $0.02$--$0.05$ m/s (slowest, with programmed delay)
\end{itemize}

\section{ECG Generation: The Dipole Model}

\subsection{Cardiac Electrical Dipole}

At any instant, the heart has a net electrical dipole moment $\mathbf{p}(t)$ pointing from negative to positive charge:

$$\mathbf{p}(t) = \sum_{\text{all myocytes}} q_i(t) \mathbf{r}_i$$

\noindent where $q_i(t)$ is the charge state of myocyte $i$ at position $\mathbf{r}_i$.

During the AP:
\begin{itemize}
    \item \textbf{Depolarization phase 0}: Wavefront of depolarization creates a dipole pointing from rest (negative) to depolarized (positive)
    \item \textbf{Plateau phase 2}: Little dipole moment (most of heart is depolarized)
    \item \textbf{Repolarization phase 3}: Wavefront of repolarization creates a dipole pointing from depolarized (negative) to rest (positive)
\end{itemize}

\subsection{12-Lead ECG and Lead Vectors}

Each of the 12 ECG leads measures the projection of the cardiac dipole onto a fixed spatial direction called the \textbf{lead vector}.

The 12 leads comprise:

\begin{itemize}
    \item \textbf{Limb leads}: I, II, III (measure frontal plane: left-right and inferior-superior)
    \item \textbf{Augmented leads}: aVR, aVL, aVF (also frontal plane, normalized)
    \item \textbf{Precordial leads}: V1--V6 (measure transverse plane: anterior-posterior and left-right)
\end{itemize}

Standard lead vectors (in mV per mV of dipole component) are fixed by convention and embedded in the simulator.

\subsection{Gaussian Dipole Model}

In this simulator, each anatomical node (e.g., RV, LV lateral) contributes to the ECG as a \textbf{time-localized dipole}:

\begin{equation}
\label{eq:dipole_gaussian}
\mathbf{p}_{\text{node}}(t) = \mathbf{d}_{\text{node}} \left[ g_{\text{depol}}(t-t_{\text{act}}) + A_{\text{repol}} \, g_{\text{repol}}(t-t_{\text{act}}) \right]
\end{equation}

where:

\begin{itemize}
    \item $\mathbf{d}_{\text{node}}$ is the \textbf{dipole vector} (anatomy-specific, units: mV-m)
    \item $t_{\text{act}}$ is the \textbf{activation time} (when the node first crossed threshold)
    \item $g_{\text{depol}}(t)$ is the \textbf{depolarization Gaussian}: $g_d(t) = \exp\left(-\frac{t^2}{2\sigma_d^2}\right) \Theta(t)$ (active for $t > 0$)
    \item $\sigma_d \approx 12.7$ ms (controls QRS width; Phase 1/2 value)
    \item $g_{\text{repol}}(t)$ is the \textbf{repolarization Gaussian}: $g_r(t) = \exp\left(-\frac{(t-\tau_r)^2}{2\sigma_r^2}\right)$ (peaks at repolarization delay $\tau_r \approx 280$ ms)
    \item $A_{\text{repol}} \approx -0.3$ (T-wave is opposite polarity to QRS)
\end{itemize}

The total \ecg voltage in lead $k$ is the sum over all nodes:

\begin{equation}
\label{eq:ecg_lead}
\text{ECG}_k(t) = \sum_{\text{nodes}} \mathbf{L}_k \cdot \mathbf{p}_{\text{node}}(t)
\end{equation}

where $\mathbf{L}_k$ is the \textbf{lead vector} for lead $k$ (dimensionless, normalized).

\section{\fhn Model (Phase 3)}

\subsection{Biological Motivation}

The \fhn model is a two-variable simplified model of the Hodgkin-Huxley equations that captures the essential fast-slow dynamics of cardiac action potentials without explicitly modeling individual ion channels.

\subsection{The FitzHugh-Nagumo Equations}

In normalized units, the 1-D cable equations are:

\begin{align}
\frac{\partial v}{\partial t} &= \frac{D}{\tau} \frac{\partial^2 v}{\partial x^2} + \frac{1}{\tau} \left( v - \frac{v^3}{3} - w + I_{\text{ext}} \right) \label{eq:fhn_v}\\
\frac{\partial w}{\partial t} &= \frac{\varepsilon}{\tau} \left( v + a - b \cdot w \right) \label{eq:fhn_w}
\end{align}

where:

\begin{itemize}
    \item $v(x,t)$ is the \textbf{transmembrane voltage} (normalized: $-1.2 \le v \le 2$; unphysiological units for mathematical simplicity)
    \item $w(x,t)$ is the \textbf{slow recovery variable} (models inactivation of inward current and activation of outward current)
    \item $x$ is the position along the cable
    \item $D$ is the \textbf{diffusion coefficient} (cell-to-cell coupling)
    \item $\tau$ is the \textbf{physical time scale} (all times scale inversely with $\tau$)
    \item $\varepsilon$ is the \textbf{time scale separation factor} (small $\varepsilon \Rightarrow$ slow $w$ recovery)
    \item $a, b$ are constants ($a \approx 0.7$, $b \approx 0.8$)
    \item $I_{\text{ext}}$ is external stimulus current (applied at specific locations during appropriate times)
\end{itemize}

\subsubsection{Fast-Slow Dynamics}

The key insight is that $v$ evolves on the \textbf{fast timescale} $\sim \tau$, while $w$ evolves on the \textbf{slow timescale} $\sim \tau/\varepsilon$:

\begin{itemize}
    \item \textbf{Depolarization (Phase 0)}: $v$ rapidly increases as the depolarizing current dominates
    \item \textbf{Plateau (Phase 2)}: $v$ is high, $w$ slowly increases, gradually opposing the depolarizing current
    \item \textbf{Repolarization (Phase 3)}: Once $w$ is large enough, it repolarizes $v$ back toward rest
    \item \textbf{Refractory period}: $w$ remains large, suppressing $v$ from re-depolarizing. As $w$ slowly decays, $v$ can eventually respond again
\end{itemize}

\subsubsection{Action Potential Duration}

The AP duration (APD) is determined by the recovery timescale:

$$\text{APD} \sim \frac{\tau}{\varepsilon \cdot b}$$

For \phaseIII with $\tau = 15$ ms, $\varepsilon = 0.08$, $b = 0.8$:

$$\tau_w = \frac{0.015}{0.08 \times 0.8} = 0.234 \text{ s} = 234 \text{ ms}$$

This recovery time constant drives the refractoriness: ventricular cells take $\approx 234$ ms to fully recover. At 70 bpm (cycle length = 857 ms), a cell fires, then 857 ms later is stimulated again when it has recovered $857 / 234 \approx 91\%$ of the way. Multi-beat ECGs therefore emerge naturally.

\subsection{Discretization: Monodomain Cable}

The cable is discretized into $N_{\text{cells}}$ compartments. The Laplacian $\frac{\partial^2 v}{\partial x^2}$ is approximated by finite differences:

$$\left.\frac{\partial^2 v}{\partial x^2}\right|_i \approx \frac{v_{i-1} - 2v_i + v_{i+1}}{\Delta x^2}$$

with no-flux (Neumann) boundary conditions at the ends:

$$v_0' = v_1 - v_0, \quad v_{N-1}' = v_{N-2} - v_{N-1}$$

\subsection{Numerical Integration}

The discretized system is integrated using explicit Euler with timestep $\Delta t = 0.1$ ms:

\begin{align}
v_i^{n+1} &= v_i^n + \Delta t \left[ \frac{D}{\Delta x^2 \tau}(v_{i-1}^n - 2v_i^n + v_{i+1}^n) + \frac{1}{\tau}(v_i^n - \frac{(v_i^n)^3}{3} - w_i^n + I_i^n) \right]\\
w_i^{n+1} &= w_i^n + \Delta t \frac{\varepsilon}{\tau}(v_i^n + a - b \cdot w_i^n)
\end{align}

The CFL (Courant-Friedrichs-Lewy) stability condition requires:

$$D \cdot \frac{\Delta t}{\tau \Delta x^2} \le 0.5$$

For the current values ($D_{\text{max}} = 72$, $\Delta t = 0.1$ ms, $\tau = 15$ ms, $\Delta x = 1$):

$$\text{CFL} = 72 \times \frac{0.0001}{0.015 \times 1^2} = 0.48 < 0.5 \quad \checkmark$$

\section{Conduction Graph Model (Phase 1/2)}

\subsection{Graph Structure}

The heart is modeled as a directed acyclic graph (DAG) with fixed \textbf{nodes} (anatomical regions) and fixed \textbf{edges} (conduction pathways):

\begin{equation}
\text{Graph} = (V, E)
\end{equation}

where:

\begin{itemize}
    \item $V = \{\text{SA, RA, LA, AV node, His, LBB, RBB, SEPT, LV, RV, etc.}\}$ (13 nodes in total)
    \item $E = \{(u,v): \text{activation can propagate from } u \text{ to } v\}$
\end{itemize}

\subsection{Activation Propagation}

At each heartbeat, the \sa node fires, triggering a \textbf{breadth-first search (BFS)} wave propagation:

\begin{algorithm}
    \caption{Activation Propagation (BFS)}
    \begin{enumerate}
        \item Queue $\leftarrow$ [SA node]
        \item Activation times $t_{\text{act}} \leftarrow$ {}
        \item While Queue is not empty:
        \begin{enumerate}
            \item $u \leftarrow$ Dequeue
            \item For each neighbor $v$ of $u$ in Graph:
            \begin{enumerate}
                \item If $v$ not yet activated:
                \item Delay $\Delta t_{\text{edge}} \leftarrow$ conduction delay on edge $(u,v)$
                \item Refractory check: if $(t_{\text{current}} - t_{\text{act}}[v]) \le$ refractory period of $v$, skip
                \item Otherwise: $t_{\text{act}}[v] \leftarrow t_{\text{act}}[u] + \Delta t_{\text{edge}}$
                \item Enqueue $v$
            \end{enumerate}
        \end{enumerate}
    \end{enumerate}
\end{algorithm}

Each node has a \textbf{refractory period} (hardcoded, not emergent):

\begin{itemize}
    \item \sa, atrial nodes: $300$ ms (atrial refractoriness)
    \item \av node: $300$ ms
    \item Ventricular nodes: $250$ ms
\end{itemize}

The \textbf{conduction delay} on each edge is also fixed:

\begin{itemize}
    \item SA to RA: $0$ ms (instantaneous)
    \item RA to LA: $100$ ms
    \item LA to AV node: $50$ ms
    \item AV node programmed delay: $100$ ms (controlled by parameter)
    \item AV to His: $0$ ms
    \item His to ventricular nodes: node-dependent ($0$--$50$ ms)
\end{itemize}

\subsection{Pathological Modifications}

Pathologies are modeled by \textbf{dynamically modifying} the graph structure or conduction parameters:

\begin{itemize}
    \item \textbf{AV Block III°}: Remove all edges from LA to AV node
    \item \textbf{RBBB}: Increase RBB conduction delay by $50$ ms
    \item \textbf{Sinus bradycardia}: Increase SA cycle length
    \item \textbf{Atrial fibrillation}: Replace BFS with chaotic re-entrant activation
\end{itemize}

% ================================================================
\chapter{Implementation in Python: Core Modules and Functions}
% ================================================================

This chapter describes the Python implementation of all models in the simulator.

\section{Project Directory Structure}

\begin{lstlisting}
cardiac_sim/
+-- core/                    # Core computation layer
|   +-- cell_models/        # Ionic models (FHN, Aliev-Panfilov, etc.)
|   +-- conduction/         # Conduction graph (Phase 1/2)
|   +-- tissue/             # Cable model (Phase 3)
|   +-- ecg/                # ECG generation (lead fields, dipoles)
|   +-- interfaces.py       # Abstract base classes
|   +-- parameter_model.py  # Parameter definitions and storage
|   └-- simulation_worker.py # Worker thread for simulation loop
+-- simulation/             # High-level simulation engines
|   +-- engine.py          # Abstract SimulationEngine
|   +-- graph_engine.py    # Phase 1/2 implementation
|   +-- cable_engine.py    # Phase 3 implementation
|   └-- solver.py          # Numerical solvers (stub)
+-- gui/                    # PyQt6 GUI layer
|   +-- main_window.py     # Main window, toolbar, layout
|   +-- ecg_display.py     # ECG waveform plotter
|   +-- parameter_panel.py # Parameters tree widget
|   └-- pathology_panel.py # Pathology selector
+-- plugins/               # Pathology plug-in system
|   +-- av_blocks.py       # AV block pathologies
|   +-- bundle_blocks.py   # Bundle branch blocks
|   +-- atrial.py          # Atrial pathologies
|   +-- ventricular.py     # Ventricular pathologies
|   └-- sinus_rhythms.py   # SA node disorders
└-- main.py               # Application entry point
\end{lstlisting}

\section{Core Interfaces and Abstract Classes}

\subsection{SimulationEngine Abstract Base Class}

File: \texttt{cardiac\_sim/core/interfaces.py}

\begin{lstlisting}
class SimulationEngine(ABC):
    @abstractmethod
    def initialize(self, params: SimulationParameters) -> None:
        """Initialize engine with parameters."""
        pass
    
    @abstractmethod
    def step(self, dt_s: float) -> ECGSample:
        """Advance by dt_s seconds, return ECG sample."""
        pass
    
    @abstractmethod
    def apply_pathology(self, plugin: AbstractPathologyPlugin) -> None:
        """Apply a pathological modification."""
        pass
    
    @abstractmethod
    def get_state(self) -> SimulationState:
        """Get current simulation state (HR, timing, etc.)."""
        pass
\end{lstlisting}

All engines (graph, cable) inherit from this abstract class and must implement all methods.

\subsection{ECGSample and SimulationState Data Classes}

\begin{lstlisting}
@dataclass
class ECGSample:
    timestamp: float      # Current time in seconds
    leads: np.ndarray    # Shape (12,) with voltages in mV
    
@dataclass
class SimulationState:
    time: float
    heart_rate: float    # BPM
    is_running: bool
\end{lstlisting}

\section{\phaseI Implementation: ConductionGraphEngine}

File: \texttt{cardiac\_sim/simulation/graph\_engine.py} (≈ 400 lines)

\subsection{Class Overview}

\begin{lstlisting}
class ConductionGraphEngine(SimulationEngine):
    def __init__(self):
        self._params: SimulationParameters = None
        self._nodes: dict[str, ECGNode] = {}      # node name → node object
        self._graph: dict[str, list[str]] = {}    # adjacency list
        self._beat: BeatRecord = None             # current beat
        self._time: float = 0.0
        self._next_fire: float = 0.0              # next SA fire time
    
    def initialize(self, params: SimulationParameters) -> None:
        self._params = params
        self._nodes = _build_nodes()              # from conduction/graph.py
        self._graph = build_physiological_graph()
        self._reset()
    
    def step(self, dt_s: float) -> ECGSample:
        self._time += dt_s
        
        # Check if SA should fire (based on HR)
        if self._time >= self._next_fire:
            self._fire_sa()
        
        # Compute ECG from active beat
        ecg = self._compute_ecg(self._time)
        return ECGSample(timestamp=self._time, leads=ecg)
    
    def _fire_sa(self) -> None:
        # BFS propagation from SA node
        self._beat = self._run_bfs()
        self._next_fire += self._params.sa_node.cycle_length_ms / 1000.0
    
    def _run_bfs(self) -> BeatRecord:
        # Breadth-first search to compute activation times
        # Returns a BeatRecord with all node activation times
        ...
    
    def _compute_ecg(self, t_s: float) -> np.ndarray:
        # Sum Gaussian dipoles from all active nodes
        ...
\end{lstlisting}

\subsection{Key Functions}

\subsubsection{\texttt{\_build\_nodes} and \texttt{build\_physiological\_graph}}

File: \texttt{cardiac\_sim/core/conduction/graph.py}

These functions construct the fixed conduction graph and node parameters:

\begin{lstlisting}
def _build_nodes() -> dict[str, ECGNode]:
    """Return dict of anatomical nodes with ECG dipole parameters."""
    nodes = {
        "SA_NODE": ECGNode(
            dipole_vector=[-1.0, 0.5, 0.2],  # mV-m, normalized
            depol_duration=0.010,              # 10 ms Gaussian width
            repol_delay=0.280,                 # T-wave peaks at 280 ms
            repol_duration=0.030,
            repol_amplitude=-0.3
        ),
        "RA": ECGNode(...),  # Right atrial myocardium
        "LA": ECGNode(...),  # Left atrial myocardium
        # ... (10 more nodes)
    }
    return nodes

def build_physiological_graph() -> dict[str, list[str]]:
    """Return adjacency list of conduction pathways."""
    return {
        "SA_NODE": ["RA"],
        "RA": ["LA", "AV_NODE"],
        "LA": ["AV_NODE"],
        "AV_NODE": ["HIS"],
        "HIS": ["LBB", "RBB", "SEPT_EARLY"],
        # ... (more edges)
    }
\end{lstlisting}

\subsubsection{\texttt{\_compute\_ecg}}

Sums Gaussian dipoles (Eq.~\ref{eq:dipole_gaussian}) over all nodes:

\begin{lstlisting}
def _compute_ecg(self, t_s: float) -> np.ndarray:
    dipole_total = np.zeros(3)  # 3D dipole vector
    
    if self._beat is not None:
        for node_name, t_act in self._beat.activation_times.items():
            node = self._nodes[node_name]
            
            # Depolarization Gaussian
            g_depol = np.exp(-(t_s - t_act)**2 / (2 * node.depol_duration**2))
            if t_s < t_act:
                g_depol = 0
            
            # Repolarization Gaussian (peaks at t_act + repol_delay)
            t_repol_peak = t_act + node.repol_delay
            g_repol = np.exp(-(t_s - t_repol_peak)**2 / (2 * node.repol_duration**2))
            
            # Total dipole contribution
            dipole = node.dipole_vector * (g_depol - node.repol_amplitude * g_repol)
            dipole_total += dipole
    
    # Project onto 12 lead vectors
    ecg_12lead = self._lead_field.project(dipole_total)  # shape (12,)
    return ecg_12lead
\end{lstlisting}

\section{\phaseIII Implementation: ConductionCableEngine}

File: \texttt{cardiac\_sim/simulation/cable\_engine.py} (≈ 450 lines)

\subsection{Class Overview}

\begin{lstlisting}
class ConductionCableEngine(SimulationEngine):
    def __init__(self):
        self._cable: ConductionCable1D = None  # 1-D FHN cable
        self._nodes: dict[str, ECGNode] = {}   # modified depol_duration
        self._params: SimulationParameters = None
        self._time: float = 0.0
        self._next_sa_fire: float = 0.0
        self._active_beats: list[BeatRecord] = []
        self._beat_in_progress: bool = False
        self._current_beat_nodes: dict[str, float] = {}
    
    def initialize(self, params: SimulationParameters) -> None:
        self._params = params
        self._cable = ConductionCable1D(params)
        self._nodes = _build_cable_nodes()  # adjusted Gaussians
        self._time = 0.0
    
    def step(self, dt_s: float) -> ECGSample:
        self._time += dt_s
        
        # SA pacemaker
        if self._time >= self._next_sa_fire:
            self._fire_sa()
        
        # Always advance cable (to keep internal clock synchronized)
        newly = self._cable.advance(dt_s)
        
        # If beat in progress, collect nodes as they activate
        if self._beat_in_progress:
            for node, t_act in newly.items():
                if node not in self._current_beat_nodes:
                    self._current_beat_nodes[node] = t_act
                    # KEY FIX: update active beat immediately so ECG
                    # renders the QRS at activation time, not at commit
                    if self._active_beats:
                        self._active_beats[-1].activation_times[node] = t_act
            
            # Commit beat when done
            ... (beat completion logic)
        
        # Compute ECG from all active beats
        ecg = self._compute_ecg(self._time)
        return ECGSample(timestamp=self._time, leads=ecg)
\end{lstlisting}

\subsection{Cable Propagation: ConductionCable1D}

File: \texttt{cardiac\_sim/core/tissue/cable\_1d.py} (≈ 500 lines)

\subsubsection{Class Overview and Fixed Parameters}

\begin{lstlisting}
# Fixed FHN model parameters (Phase 3-D calibrated)
_A   = 0.700     # FHN parameter a
_B   = 0.800     # FHN parameter b
_EPS = 0.080     # recovery rate
_TAU = 0.015     # time scale (15 ms for correct PR/QRS timing)

_DT_INT = 0.0001 # internal micro-step (0.1 ms)

# Cable structure
_NA = 8          # atrial cells (SA + 7)
_NV = 10         # ventricular cells (His + LBB(3) + LV(6))

# Diffusion coefficients
_D_ATR = np.array([0.0, 2.88, 2.88, 2.88, 2.88, 2.88, 2.88, 2.88])
_D_VEN = np.array([0.0, 72.0, 72.0, 72.0, 3.0, 3.0, 3.0, 3.0, 3.0, 3.0])

# Stimulus
_STIM_AMP = 1.5  # suprathreshold current
_STIM_DUR = 0.010 # 10 ms

class ConductionCable1D:
    def __init__(self, params: SimulationParameters):
        self._sa: np.ndarray = np.empty((_NA, 2))  # (cells, [v, w])
        self._sv: np.ndarray = np.empty((_NV, 2))
        self._t: float = 0.0
        self._check_cfl()
    
    def new_beat(self, sa_fire_time: float, params: SimulationParameters) -> None:
        """Reset for a new beat at sa_fire_time."""
        self._t = sa_fire_time
        self._stim_end_a = sa_fire_time + _STIM_DUR
        self._av_pending_at = math.inf
        # Reset activation trackers
        self._t_act_a[:] = math.inf
        self._t_act_v[:] = math.inf
        self._t_act_a[0] = sa_fire_time
    
    def advance(self, dt_s: float) -> dict[str, float]:
        """Integrate by dt_s seconds, return newly activated nodes."""
        n_micro = max(1, round(dt_s / _DT_INT))
        dt_micro = dt_s / n_micro
        newly: dict[str, float] = {}
        
        for _ in range(n_micro):
            self._t += dt_micro
            self._step_atrial(dt_micro)
            self._step_ventricular(dt_micro)
            
            # Check for threshold crossings
            for node, (i0, i1) in _ATR_SEGS.items():
                # (threshold detection logic)
                ...
            
            for node, (i0, i1) in _VEN_SEGS.items():
                # (threshold detection logic)
                ...
        
        return newly
    
    def _step_atrial(self, dt: float) -> None:
        """One FHN Euler step on atrial cable."""
        V = np.ascontiguousarray(self._sa[:, 0])
        W = np.ascontiguousarray(self._sa[:, 1])
        I = np.zeros(_NA)
        if self._t < self._stim_end_a:
            I[0] = _STIM_AMP
        V_new, W_new = _fhn_cable_step(V, W, _D_ATR, I, dt, _TAU, _A, _B, _EPS)
        self._sa[:, 0] = V_new
        self._sa[:, 1] = W_new
\end{lstlisting}

\subsubsection{Numerical Integration: \texttt{\_fhn\_cable\_step}}

File: \texttt{cardiac\_sim/core/tissue/cable\_1d.py}

\begin{lstlisting}
def _fhn_cable_step(V, W, D, I_ext, dt, tau_s, a, b, eps):
    """Single explicit Euler step for FHN cable.
    
    Implements equations (1) and (2) with finite-difference Laplacian.
    """
    # Compute Laplacian with no-flux BCs
    lap = np.empty_like(V)
    lap[1:-1] = V[:-2] - 2.0 * V[1:-1] + V[2:]
    lap[0]    = V[1] - V[0]
    lap[-1]   = V[-2] - V[-1]
    
    inv_tau = 1.0 / tau_s
    dv = (D * lap + V - V**3 / 3.0 - W + I_ext) * inv_tau
    dw = eps * (V + a - b * W) * inv_tau
    
    return V + dt * dv, W + dt * dw
\end{lstlisting}

This function is Numba-JIT compilable (if compatible NumPy version) for 10-20$\times$ speedup.

\end{document}
